% Source: https://artofproblemsolving.com/community/c3h2650413_sigma_mathematics_competition_smc_problems
% Stable identity: aops-topic-2650413-q03
We need $n>2$; put $m=n-2\ge1$, so the expression is $\dfrac{m+8}{\sqrt m}$. If it equals an integer $q$, then $\sqrt m=\dfrac{m+8}{q}$ is rational, so $m=k^2$ for a positive integer $k$. Then $\dfrac{k^2+8}{k}=k+\dfrac8k$ is an integer exactly when $k\mid8$, i.e. $k\in\{1,2,4,8\}$. This gives $n=k^2+2\in\{3,6,18,66\}$ (values $9,6,6,9$), with sum $\boxed{93}$.
